\chapter{Az AES titkosítás és különösen a GCM mód}

Ebben a fejezetben összefoglalom az AES blokk-kód alapötletét, majd a projektben használt
\textbf{CTR} és \textbf{GCM} mód lényegét. A hangsúly a GCM-en van, mert ez egyszerre ad
\emph{titkosságot} és \emph{integritást/hitelességet} (AEAD: Authenticated Encryption with Associated Data).

\section{AES rövid áttekintés}

Az \textbf{AES (Advanced Encryption Standard)} egy 128 bites blokk-kód, amely
\textbf{128/192/256 bites} kulcsokkal működik. A titkosítás több \emph{round}-ból áll
(10/12/14 round a kulcsmérettől függően), és minden round tipikusan négy klasszikus lépést tartalmaz:

\begin{itemize}
  \item \textbf{SubBytes}: nemlineáris S-box helyettesítés bájtonként,
  \item \textbf{ShiftRows}: sor-eltolás a 4x4-es state mátrixban,
  \item \textbf{MixColumns}: lineáris keverés oszloponként (GF($2^8$) felett),
  \item \textbf{AddRoundKey}: round-kulcs XOR-olása a state-tel.
\end{itemize}

A blokk-kód önmagában csak fix méretű (128 bit) blokkokra értelmezett, ezért gyakorlati használathoz
üzemmódra van szükség.

\section{CTR mód (Counter mode)}

A \textbf{CTR} (counter) mód stream-szerű titkosítást ad: egy számlálóval (counter) generálok
egy \emph{kulcsfolyamot} (\emph{keystream}), amit XOR-olok a plaintexttel.

\begin{align*}
  S_i &= AES_K(IV \| counter_i) \\\\
  C_i &= P_i \oplus S_i
\end{align*}

Fontos tulajdonságok:
\begin{itemize}
  \item A blokkok \emph{függetlenek}: $S_i$ számolható párhuzamosan $\rightarrow$ CTR kiválóan párhuzamosítható GPU-n is.
  \item Az IV/counter párnak \emph{egyedinek} kell lennie ugyanazzal a kulccsal (különben kulcsfolyam-újrafelhasználás lép fel).
\end{itemize}

\section{GCM mód (Galois/Counter Mode) -- miért több, mint CTR?}

A CTR önmagában \textbf{nem} ad integritást: a ciphertext módosítható, és a decrypt után
ellenőrzés nélkül hibás plaintext jön ki. A \textbf{GCM} ezt kiegészíti egy
\textbf{hitelesítő tag} (authentication tag) számításával.

A GCM két fő része:
\begin{enumerate}
  \item \textbf{Titkosítás CTR-rel} (mint fent),
  \item \textbf{GHASH} -- egy polinomiális MAC-szerű összegzés GF($2^{128}$) felett, amely az AAD-t és a ciphertextet összegzi.
\end{enumerate}

\subsection{A GHASH alapötlete}

A GCM kiszámít egy hash-szubrutin kulcsot:
\[
H = AES_K(0^{128})
\]
Majd a \textbf{GHASH} az AAD és a ciphertext blokkokat egy GF($2^{128}$) mezőben értelmezve
szorozza/összegzi $H$-val. Ebből jön ki egy 128 bites érték, amit a végén még CTR-rel
``letitkosít'' a tag előállításához.

\subsection{GCM adatfolyam (ábra)}

Az alábbi ábra a projektben is használt tipikus (12 bájtos) IV-es GCM folyamatot mutatja.

\begin{figure}[h!]
  \centering
  \begin{tikzpicture}[
    node distance=10mm,
    box/.style={draw, rectangle, rounded corners=2pt, align=center, minimum width=3.8cm, minimum height=0.9cm},
    small/.style={draw, rectangle, rounded corners=2pt, align=center, minimum width=2.9cm, minimum height=0.8cm},
    arrow/.style={-Latex, thick}
  ]
    \node[box] (pt) {Plaintext $P$};
    \node[box, right=18mm of pt] (ctr) {AES-CTR\\(keystream + XOR)};
    \node[box, right=18mm of ctr] (ct) {Ciphertext $C$};

    \node[small, below=14mm of ctr] (aad) {AAD};
    \node[box, below=14mm of ct] (ghash) {GHASH\\GF($2^{128}$)};
    \node[box, left=18mm of ghash] (tag) {Tag $T$};

    \node[small, below=14mm of tag] (j0) {$J_0$ (IV-ből)};
    \node[small, below=14mm of ghash] (h) {$H = AES_K(0^{128})$};

    \draw[arrow] (pt) -- (ctr);
    \draw[arrow] (ctr) -- (ct);

    \draw[arrow] (aad) -- (ghash);
    \draw[arrow] (ct) -- (ghash);

    \draw[arrow] (h) -- (ghash);
    \draw[arrow] (j0) -- (tag);
    \draw[arrow] (ghash) -- (tag);

    % annotate: tag calc
    \node[below=2mm of tag, align=center] (tagformula) {\scriptsize $T = AES_K(J_0) \oplus GHASH_H(AAD, C)$};
  \end{tikzpicture}
  \caption{AES-GCM magas szintű folyamat: CTR titkosítás + GHASH + tag előállítás.}
\end{figure}

\subsection{Dekódolás / ellenőrzés}

Visszafejtéskor először a megkapott $AAD$ és $C$ alapján újraszámolom a tag-et,
és \textbf{csak akkor fogadom el a plaintextet}, ha a tag egyezik. Ez a lényegi különbség
a sima CTR-hez képest: a GCM \emph{hitelesített titkosítás}.

\section{Miért érdekes ez teljesítmény szempontból?}

\begin{itemize}
  \item CTR-ben a blokkok függetlensége miatt a GPU-s párhuzamosítás nagyon természetes.
  \item GCM-ben a CTR rész ugyanúgy jól skálázódik, de a GHASH egy része és a tag kezelés
        \emph{szekvenciális komponenseket} is bevisz a folyamatba. Emiatt a gyorsítás mértéke
        általában kisebb, mint a ``tiszta'' CTR-ben.
  \item Emellett a GPU-s útvonalnál a host--device adatmozgatás (\emph{memcpy/PCIe}) fix költsége
        kisebb adatmennyiségnél könnyen dominálhat.
\end{itemize}

\cleardoublepage
